\section{Processo organizacional — AS-IS consolidado}
O processo atual inicia-se com um emissor identificado que elabora ou solicita a divulgação de uma informação. O conteúdo pode ser encaminhado informalmente, por diferentes canais ou por avaliação de direção e coordenação. Após a publicação, não existe garantia homogênea de que os destinatários receberam, visualizaram ou compreenderam a informação.

\begin{enumerate}
\item Identificar o emissor e a necessidade de comunicação.
\item Elaborar o comunicado.
\item Selecionar canais, público-alvo e, eventualmente, responsáveis pela aprovação.
\item Solicitar ou realizar avaliação quando necessário.
\item Publicar ou distribuir por canais oficiais ou informais.
\item Comunicar o resultado ao emissor.
\item Receber, consultar ou encaminhar o aviso ao público interessado.
\end{enumerate}

O AS-IS apresenta fragmentação, falta de rastreabilidade, decisões editoriais inconsistentes e ausência de confirmação de recebimento. Os principais problemas são perdas de prazos, aumento de ruído e dificuldade de segmentação das informações.
